\chapter{Diseases of the Pancreas}
\label{ch:pancreas}
The pancreas is a retroperitoneal organ with both exocrine (digestive enzyme production) and endocrine (insulin and glucagon secretion) functions. The exocrine pancreas produces approximately 1.5 liters of pancreatic juice daily, containing bicarbonate and digestive enzymes (lipase, amylase, trypsinogen, chymotrypsinogen). The endocrine pancreas consists of the islets of Langerhans, containing beta cells (insulin), alpha cells (glucagon), delta cells (somatostatin), and PP cells (pancreatic polypeptide). Diseases of the pancreas include acute and chronic pancreatitis, diabetes mellitus, and pancreatic cancer.

%-------------------------------------------------------------
\section{Diabetes Mellitus}
%-------------------------------------------------------------

%-------------------------------------------------------------

%-------------------------------------------------------------

%-------------------------------------------------------------

Diabetes mellitus is a chronic metabolic disorder characterized by persistent hyperglycemia resulting from defects in insulin secretion, insulin action, or both. Absolute or relative insulin deficiency disrupts systemic glucose utilization and lipid metabolism. Uncontrolled hyperglycemia drives microvascular (retinopathy, nephropathy, neuropathy) and macrovascular (coronary artery disease, stroke) complications requiring lifelong glycemic control.

%-------------------------------------------------------------

%-------------------------------------------------------------

%-------------------------------------------------------------

%-------------------------------------------------------------

\label{sec:Diabetes Mellitus}
%-------------------------------------------------------------%-------------------------------------------------------------

\subsection{Diabetic Ketoacidosis}
\label{sec:Diabetic Ketoacidosis}
Diabetic ketoacidosis (DKA) is a life-threatening metabolic emergency characterized by hyperglycemia (glucose >250 mg/dL), metabolic acidosis (pH <7.3, bicarbonate <18 mEq/L), and ketosis (elevated ketones in blood and urine). It occurs predominantly in T1DM but can occur in T2DM under significant physiological stress. Precipitants include insulin non-compliance, infections, new-onset diabetes, myocardial infarction, and medications. The condition results from absolute or relative insulin deficiency counteracted by counter-regulatory hormones (glucagon, cortisol, catecholamines), leading to lipolysis, ketone body production (beta-hydroxybutyrate, acetoacetate), and osmotic diuresis causing dehydration and electrolyte depletion. Clinical features include polyuria, polydipsia, nausea, vomiting, abdominal pain, Kussmaul respirations (deep, rapid breathing), and a fruity breath odor. Diagnosis is based on blood glucose, arterial blood gas, serum ketones, and anion gap calculation. Management involves aggressive IV fluid resuscitation, continuous IV insulin infusion, and careful potassium replacement. Prognosis is favorable with prompt treatment; mortality is low in well-managed cases.

\subsection{Gestational Diabetes Mellitus}
\label{sec:Gestational Diabetes Mellitus}
Gestational diabetes mellitus (GDM) is glucose intolerance first recognized during pregnancy, affecting 2--10\% of pregnancies. The condition results from insufficient insulin secretion to overcome pregnancy-related insulin resistance mediated by placental hormones (human placental lactogen, cortisol, progesterone). Risk factors include obesity, advanced maternal age, family history of diabetes, prior GDM, and polycystic ovary syndrome. Clinical features are typically asymptomatic; screening is performed at 24--28 weeks gestation using the 75-g oral glucose tolerance test. Complications for the mother include preeclampsia and increased cesarean delivery rate; fetal/neonatal complications include macrosomia, shoulder dystocia, neonatal hypoglycemia, and long-term risk of obesity and T2DM. Diagnosis is based on the oral glucose tolerance test. Management includes medical nutrition therapy, exercise, and insulin therapy when glycemic targets are not met with lifestyle measures; metformin may be used in select cases. Prognosis is generally favorable; GDM usually resolves postpartum, but women with GDM have a 50\% lifetime risk of developing T2DM and should undergo periodic screening.

\subsection{Hyperosmolar Hyperglycemic State}
\label{sec:Hyperosmolar Hyperglycemic State}
Hyperosmolar hyperglycemic state (HHS) is a metabolic emergency primarily affecting T2DM patients, characterized by extreme hyperglycemia (glucose >600 mg/dL), profound dehydration, hyperosmolality (serum osmolality >320 mOsm/kg), and minimal or absent ketosis. The condition results from sufficient residual insulin to prevent ketogenesis but not to control glucose. Clinical features include severe dehydration, altered mental status, and neurological deficits (seizures, focal deficits). Diagnosis is based on blood glucose, serum osmolality, and absence of significant ketosis. Management involves aggressive fluid resuscitation (initially normal saline), insulin infusion (at lower doses than DKA), and electrolyte replacement. Prognosis is more guarded than DKA; mortality is higher (up to 20\%).

\subsection{Hypoglycemia}
\label{sec:Hypoglycemia}
Hypoglycemia is a blood glucose level below 70 mg/dL, with clinically significant hypoglycemia occurring below 54 mg/dL (Level 2) and severe hypoglycemia (Level 3) being any event requiring external assistance. It is the most common acute complication of diabetes treatment (insulin and sulfonylureas). The condition results from an imbalance between glucose-lowering therapy and glucose supply. Clinical features include adrenergic symptoms (tremor, palpitations, sweating, anxiety) and neuroglycopenic symptoms (confusion, blurred vision, seizures, loss of consciousness). Diagnosis is based on the Whipple triad: symptoms of hypoglycemia, low measured glucose, and resolution of symptoms with glucose correction. Management of mild-moderate hypoglycemia is with rapid-acting oral glucose (15--20 g); severe hypoglycemia requires IV dextrose or glucagon injection. Prognosis is generally excellent with prompt treatment; recurrent hypoglycemia may cause hypoglycemia unawareness, requiring adjustment of insulin regimens.

\subsection{Type 1 Diabetes Mellitus}
\label{sec:Type 1 Diabetes Mellitus}
Type 1 diabetes mellitus (T1DM) is an autoimmune disease characterized by destruction of pancreatic beta cells in the islets of Langerhans, leading to absolute insulin deficiency. It accounts for 5--10\% of diabetes cases. The condition results from autoreactive T cells (with CD8+ cytotoxic T cells being primary effectors) that destroy beta cells, with autoantibodies (anti-GAD65, anti-IA-2, anti-insulin, anti-ZnT8) serving as biomarkers. Genetic susceptibility is associated with HLA-DR3/DR4 and HLA-DQ2/DQ8 haplotypes; environmental triggers (viral infections, early cow's milk exposure) may initiate autoimmunity in predisposed individuals. Clinical features include the classic triad of polyuria, polydipsia, and weight loss, with diabetic ketoacidosis (DKA) as the initial presentation in 30--40\% of children. Diagnosis requires fasting plasma glucose $\geq$126 mg/dL, random glucose $\geq$200 mg/dL with symptoms, HbA1c $\geq$6.5\%, or 2-hour OGTT glucose $\geq$200 mg/dL. Management is lifelong insulin therapy via multiple daily injections (basal-bolus regimen) or continuous subcutaneous insulin pump; continuous glucose monitoring (CGM) and insulin pump systems (closed-loop "artificial pancreas" systems) improve glycemic control. Prognosis depends on glycemic control; microvascular and macrovascular complications determine long-term outcomes. Islet transplantation and Teplizumab (anti-CD3 antibody) are emerging therapies.

\subsection{Type 2 Diabetes Mellitus}
\label{sec:Type 2 Diabetes Mellitus}
Type 2 diabetes mellitus (T2DM) is the most common form of diabetes, accounting for 90--95\% of cases. It is characterized by insulin resistance in peripheral tissues (muscle, liver, adipose tissue) combined with progressive beta-cell dysfunction and relative insulin deficiency. Risk factors include obesity (particularly visceral adiposity), physical inactivity, family history, age >45, ethnicity, history of gestational diabetes, and polycystic ovary syndrome. The pathophysiology involves the "ominous octet"---eight mechanisms contributing to hyperglycemia: decreased insulin secretion, increased glucagon secretion, increased hepatic glucose production, decreased muscle glucose uptake, increased lipolysis, decreased incretin effect, increased renal glucose reabsorption, and neurotransmitter dysfunction. Many patients are asymptomatic and diagnosed on routine screening. Clinical features of advanced disease include polyuria, polydipsia, and weight loss. Diagnosis is based on fasting plasma glucose, HbA1c, or OGTT. Complications include microvascular disease (retinopathy, nephropathy, neuropathy) and macrovascular disease (coronary artery disease, stroke, peripheral arterial disease). Management follows a stepwise approach: lifestyle modification (diet, exercise, weight loss of 5--10\%), metformin as first-line pharmacotherapy, followed by SGLT2 inhibitors, GLP-1 receptor agonists, DPP-4 inhibitors, thiazolidinediones, sulfonylureas, and insulin as needed. Prognosis depends on glycemic control and management of cardiovascular risk factors; cardiovascular benefit is demonstrated with SGLT2 inhibitors and GLP-1 receptor agonists.

%-------------------------------------------------------------
\section{Pancreatic Exocrine Insufficiency}
%-------------------------------------------------------------

%-------------------------------------------------------------

%-------------------------------------------------------------

%-------------------------------------------------------------

%-------------------------------------------------------------

%-------------------------------------------------------------

%-------------------------------------------------------------

%-------------------------------------------------------------

\label{sec:section:Pancreatic Exocrine Insufficiency}
%-------------------------------------------------------------%-------------------------------------------------------------

\subsection{Pancreatic Exocrine Insufficiency}
\label{sec:Pancreatic Exocrine Insufficiency}
Pancreatic exocrine insufficiency (PEI) is the failure of the pancreas to produce sufficient digestive enzymes, leading to maldigestion and malabsorption. Common causes include chronic pancreatitis, cystic fibrosis, pancreatic surgery (Whipple procedure, distal pancreatectomy), and advanced pancreatic cancer. The condition results from loss of pancreatic acinar cell function; lipase deficiency is the most important, leading to steatorrhea. Clinical features include fatty, foul-smelling, floating stools, weight loss, and fat-soluble vitamin deficiencies (A, D, E, K). Diagnosis is based on fecal elastase-1 level (<200 $\mu$g/g indicates moderate insufficiency, <100 $\mu$g/g severe) or 72-hour fecal fat collection. Management is pancreatic enzyme replacement therapy (PERT) with enteric-coated enzyme capsules taken with meals, with dose titrated based on symptoms and nutritional status. Prognosis is good with adequate enzyme replacement.

%-------------------------------------------------------------
\section{Pancreatic Neoplasms}
%-------------------------------------------------------------

%-------------------------------------------------------------

%-------------------------------------------------------------

%-------------------------------------------------------------

%-------------------------------------------------------------

%-------------------------------------------------------------

%-------------------------------------------------------------

%-------------------------------------------------------------

\label{sec:Pancreatic Neoplasms}
%-------------------------------------------------------------%-------------------------------------------------------------

\subsection{Pancreatic Adenocarcinoma}
\label{sec:Pancreatic Adenocarcinoma}
Pancreatic adenocarcinoma is one of the most lethal malignancies, with a 5-year survival rate of approximately 12\%. It is the fourth leading cause of cancer death worldwide. Risk factors include smoking, obesity, chronic pancreatitis, diabetes mellitus (new-onset in an older patient), family history, and BRCA2 mutations. Approximately 80\% of tumors arise in the pancreatic head. The condition results from accumulation of genetic mutations (KRAS, CDKN2A, TP53, SMAD4) driving malignant transformation. Clinical features include painless jaundice (Courvoisier's sign---palpable, non-tender gallbladder), new-onset diabetes, weight loss, epigastric pain radiating to the back, and depression. Diagnosis is based on CT abdomen with pancreatic protocol, CA 19-9 (tumor marker), and endoscopic ultrasound-guided fine needle aspiration for tissue diagnosis. Staging uses the TNM system. Management involves surgical resection (Whipple procedure for head tumors, distal pancreatectomy for body/tail tumors) as the only curative option, possible in only 15--20\% of patients; adjuvant chemotherapy (FOLFIRINOX or gemcitabine/capecitabine) improves survival. Metastatic disease is treated with systemic chemotherapy; molecular profiling may identify targets for precision therapy (BRCA mutations, MSI-high tumors). Prognosis remains poor, with a 5-year survival of approximately 12\%.

\subsection{Pancreatic Neuroendocrine Tumors}
\label{sec:Pancreatic Neuroendocrine Tumors}
Pancreatic neuroendocrine tumors (PanNETs) are neoplasms arising from the islet cells of the pancreas, accounting for 1--2\% of pancreatic neoplasms. They may be functional (hormone-secreting, 40--60\%) or non-functional. Functional PanNETs include insulinoma (most common functional type, presenting with hypoglycemia), gastrinoma (Zollinger-Ellison syndrome with recurrent peptic ulcers), glucagonoma (necrolytic migratory erythema, diabetes, weight loss), VIPoma (watery diarrhea, hypokalemia, achlorhydria), and somatostatinoma (diabetes, steatorrhea, gallstones). The condition results from neoplastic transformation of pancreatic islet cells. Clinical features depend on the hormone secreted; non-functional tumors are often discovered incidentally. Diagnosis involves hormone levels for functional tumors and imaging including CT, MRI, and somatostatin receptor scintigraphy (Octreoscan) or Ga-68 DOTATATE PET/CT. Management is surgical resection for localized disease; metastatic disease is treated with somatostatin analogs (octreotide, lanreotide), targeted therapy (sunitinib, everolimus), PRRT (peptide receptor radionuclide therapy with Lu-177 DOTATATE), and chemotherapy. Prognosis is generally better than pancreatic adenocarcinoma, depending on tumor grade and stage.

%-------------------------------------------------------------
\section{Pancreatitis}
%-------------------------------------------------------------

%-------------------------------------------------------------

%-------------------------------------------------------------

%-------------------------------------------------------------

Pancreatitis is an inflammatory disorder of the exocrine pancreas classified as acute or chronic based on clinical presentation and histological reversibility. Premature intraparenchymal activation of zymogen enzymes triggers autodigestion, local tissue necrosis, systemic inflammatory response, and potential organ failure. Etiologies predominantly include biliary tract obstruction from gallstones, heavy alcohol consumption, hypertriglyceridemia, and genetic risk factors.

%-------------------------------------------------------------

%-------------------------------------------------------------

%-------------------------------------------------------------

%-------------------------------------------------------------

\label{sec:Pancreatitis}
%-------------------------------------------------------------%-------------------------------------------------------------

\subsection{Acute Pancreatitis}
\label{sec:Acute Pancreatitis}
Acute pancreatitis is acute inflammation of the pancreas characterized by sudden onset of severe epigastric pain radiating to the back, often with nausea and vomiting. The two most common causes worldwide are gallstones (40\%) and alcohol excess (30\%); other causes include hypertriglyceridemia (triglycerides >1000 mg/dL), medications (azathioprine, valproic acid, thiazides), hypercalcemia, trauma, ERCP, and genetic factors (PRSS1, SPINK1 mutations). The condition results from premature activation of digestive enzymes within the pancreas, leading to autodigestion and inflammation. Clinical features include characteristic abdominal pain, nausea, and vomiting. Diagnosis requires two of three criteria: characteristic abdominal pain, serum amylase or lipase >3 times the upper limit of normal, and characteristic findings on imaging. Severity is assessed using APACHE II, Ranson's criteria, and modified Marshall score. Management includes aggressive IV fluid resuscitation (lactated Ringer's), pain control, early enteral nutrition, and monitoring for complications; infected pancreatic necrosis may require endoscopic, percutaneous, or surgical necrosectomy. Prognosis is generally favorable; most cases (80\%) are mild and self-limiting. Severe acute pancreatitis is characterized by organ failure (persistent for >48 hours) and/or local complications (necrosis, pseudocyst, abscess).

\subsection{Autoimmune Pancreatitis}
\label{sec:Autoimmune Pancreatitis}
Autoimmune pancreatitis (AIP) is a form of chronic pancreatitis caused by autoimmune inflammation of the pancreatic parenchyma. Type 1 AIP is associated with IgG4-related disease and involves other organs (biliary strictures, retroperitoneal fibrosis, salivary gland enlargement); Type 2 AIP is limited to the pancreas and is associated with inflammatory bowel disease. The condition results from autoimmune-mediated inflammation. Clinical features often mimic pancreatic cancer with painless jaundice, weight loss, and a pancreatic mass. Diagnosis is based on the international consensus diagnostic criteria (ICDC): imaging (sausage-shaped pancreas on CT, delayed enhancement), serology (elevated IgG4 in type 1), histology (lymphoplasmacytic infiltrate with storiform fibrosis), and response to corticosteroids. Management is with oral prednisone, with rituximab for relapsing type 1 AIP. Prognosis is favorable; unlike pancreatic cancer, AIP responds dramatically to corticosteroids.

\subsection{Autoimmune Pancreatitis (AIP)}
\label{sec:Autoimmune Pancreatitis (AIP)}
Autoimmune pancreatitis (AIP) is a distinct form of chronic pancreatitis driven by autoimmune inflammation, classified into Type 1 (IgG4-related disease, systemic involvement with dense IgG4+ plasma cell infiltration and storiform fibrosis) and Type 2 (pancreas-specific, neutrophilic phlebotomy lesions without IgG4 systemic disease). Type 1 AIP affects predominantly older males and presents with painless obstructive jaundice, weight loss, and a diffuse 'sausage-shaped' enlargement of the pancreas with a capsule-like rim on CT/MRI, frequently mimicking pancreatic ductal adenocarcinoma. Diagnosis relies on the International Consensus Diagnostic Criteria (ICDC): elevated serum IgG4 levels ($>2\times$ normal), characteristic parenchymal enlargement, main pancreatic duct narrowing, and IgG4+ plasma cells on core biopsy. AIP demonstrates dramatic responsiveness to systemic corticosteroid therapy (prednisone), avoiding unnecessary pancreaticoduodenectomy (Whipple procedure).

\subsection{Chronic Pancreatitis}
\label{sec:Chronic Pancreatitis}
Chronic pancreatitis is progressive, irreversible morphological changes of the pancreatic parenchyma leading to impairment of exocrine and endocrine function. The most common cause in Western countries is chronic alcohol use (70--80\%); other causes include hereditary pancreatitis (PRSS1 mutations), cystic fibrosis (CFTR mutations), autoimmune pancreatitis, tropical pancreatitis, and idiopathic. The condition results from repeated episodes of pancreatic injury leading to fibrosis, parenchymal atrophy, ductal strictures, and calcifications. Clinical features include chronic epigastric pain (often worse after meals), steatorrhea (exocrine insufficiency), diabetes mellitus (endocrine insufficiency), and weight loss. Diagnosis is based on CT scan showing pancreatic calcifications, ductal dilatation, and parenchymal atrophy; MRCP for ductal morphology; and functional testing (fecal elastase, secretin stimulation). Management includes pancreatic enzyme replacement therapy (PERT), pain management (analgesics, endoscopic drainage of ductal strictures, celiac plexus block), insulin for diabetes, and total pancreatectomy with islet autotransplantation for refractory cases. Prognosis is variable with progressive functional decline.

\subsection{Intraductal Papillary Mucinous Neoplasm (IPMN)}
\label{sec:Intraductal Papillary Mucinous Neoplasm (IPMN)}
Intraductal papillary mucinous neoplasm (IPMN) is a mucin-producing epithelial precursor lesion arising within the main pancreatic duct or branch ducts, characterized by intraductal papillary proliferation, mucin hypersecretion, and progressive cystic ductal ectasia. IPMNs carry variable malignant potential, progressing through low-, intermediate-, and high-grade dysplasia to invasive IPMN-associated adenocarcinoma. Classified by ductal location into main-duct (high malignant risk, $>60\%$), branch-duct (lower risk), and mixed-type. Clinical features present as recurrent acute pancreatitis (due to mucin plugging of the ampulla of Vater) or incidental pancreatic cysts on cross-sectional imaging. Mucin extruding from a gaping 'fish-mouth' ampulla on esophagogastroduodenoscopy is pathognomonic. Management is dictated by Sendai/Fukuoka guidelines: main-duct IPMNs ($>5$ mm duct) undergo surgical resection, whereas branch-duct IPMNs are managed conservatively or resected based on high-risk stigmata (solid nodule, main duct dilation, jaundice).

\subsection{Nesidioblastosis and Congenital Hyperinsulinism}
\label{sec:Nesidioblastosis and Congenital Hyperinsulinism}
Nesidioblastosis (persistent hyperinsulinemic hypoglycemia of infancy) is a rare disorder of pancreatic endocrine cell dysplasia characterized by diffuse hypertrophy, hyperplasia, and neoformation of insulin-producing pancreatic $\beta$-cells from ductal epithelium. Driven by inactivating mutations in ABCC8 or KCNJ11 genes encoding the ATP-sensitive potassium channel ($K_{\text{ATP}}$) in $\beta$-cells, the condition leads to unregulated, continuous insulin secretion unresponsive to blood glucose levels. Clinical features present in neonates or infants with severe, intractable hypoglycemia, lethargy, seizures, and risk of irreversible neurodevelopmental damage. Diagnosis is established by persistent endogenous hyperinsulinemia, elevated C-peptide, and negative sulfonylurea screening during hypoglycemia. Management involves IV glucose, diazoxide ($K_{\text{ATP}}$ channel opener), octreotide, and subtotal or near-total ($95\%$) pancreatectomy for medical treatment failure.

\subsection{Pancreatic Pseudocyst}
\label{sec:Pancreatic Pseudocyst}
Pancreatic pseudocyst is a localized fluid collection rich in pancreatic enzymes, necrotic debris, and tissue fluid enclosed by a non-epithelialized wall of fibrous and granulation tissue, occurring as a late complication ($>4$ weeks) of acute or chronic pancreatitis or pancreatic trauma. Unlike true pancreatic cysts, pseudocysts lack an epithelial lining. Driven by pancreatic ductal disruption, extravasated pancreatic juice digests surrounding tissues until contained by inflammatory adhesions of adjacent organs (stomach, duodenum, transverse mesocolon). Clinical features include persistent abdominal pain, early satiety, nausea, vomiting, and a palpable epigastric mass following an episode of pancreatitis. Complications include infection (abscess formation), gastric outlet obstruction, rupture, and pseudoaneurysm erosion into adjacent vessels. Diagnosis is established by CT or MRI showing a well-circumscribed fluid collection with a mature fibrous wall. Management of asymptomatic pseudocysts is observation; symptomatic or enlarging cysts ($>6$ cm) undergo endoscopic cystogastrostomy, percutaneous drainage, or surgical cystojejunostomy.

